\section{Introduction}

Language-model agents can now search, plan, write code, run tools, and draft
scientific text. Several systems combine these abilities into broad research
workflows \cite{lu2024aiscientist,yamada2025aiscientistv2,schmidgall2025agentlab}.
The difficult question is no longer whether an agent can generate a plausible
paper-shaped artifact. It is whether a failed result changes the next
scientific action in a way that is both substantive and auditable. Repeating
the same parameters, rewriting an explanation, or selecting a favorable
metric after seeing a test set can look iterative while adding no new
scientific evidence.

Existing evaluations illuminate important parts of this problem. Agent
benchmarks measure tool use and multi-step execution
\cite{liu2023agentbench,zhou2024webarena,xie2024osworld}. Machine-learning
benchmarks evaluate model development \cite{huang2023mlagentbench,chan2024mlebench},
and scientific-agent benchmarks measure program generation or recovery of
published computations \cite{chen2024scienceagentbench,siegel2024corebench,
starace2025paperbench}. These settings establish useful capabilities, but an
execution score alone does not reveal whether a system distinguishes a
negative scientific result from an infrastructure failure, proposes a
mechanistically different successor, freezes the successor before evaluation,
and reports the failure when the gate is missed.

Figure~\ref{fig:running-example} gives the distinction studied here. A
one-pass controller may retry or rewrite after a failed experiment. The
evidence-bound controller instead records a typed diagnosis, binds a new
hypothesis to the failed result, commits a falsification rule, and evaluates
the frozen artifact. A failed successor is not erased. It becomes the parent
evidence for another round or an explicit route change. This design treats
the research history as part of the result rather than as disposable
intermediate context.

\begin{figure*}[t]
  \centering
  \includegraphics[width=.94\textwidth]{figures/running-example.pdf}
  \Description{Two workflow paths after a failed experiment. The upper path
  retries or rewrites once. The lower path records a hash-bound diagnosis,
  proposes a new mechanism, freezes before reveal, and exports evidence and a
  report before looping to a new diagnosis.}
  \caption{Running example. A retry can change an artifact without creating a
  new scientific iteration. The evidence-bound path requires a hash-linked
  diagnosis, a changed mechanism, a pre-result freeze, deterministic
  adjudication, and a visible report.}
  \label{fig:running-example}
\end{figure*}

We operationalize this distinction in \method{}, a persistent local research
campaign. The language model proposes structured diagnoses and candidate
mechanisms, while deterministic code owns data partitioning, metric
calculation, stopping rules, and decisions. The campaign uses a shared,
versioned \vault{} as project memory. Each round produces a hypothesis,
preregistration, experiment manifest, raw metrics, validation record, failure
analysis, research report, loop report, evidence map, figures, decision, and
manuscript version. Parent and child manifests form a content-hash chain.
Unseen outcomes cannot flow back into the already frozen round.

The evaluation has two connected parts. First, two new scientific machine
learning rounds respond to an existing noise-robustness failure in equation
discovery. Round 1 tests noise-conditioned smoothing with bootstrap ensemble
sparse regression. Round 2 tests analytic spline derivatives with
cross-output group-sparse projection. Each uses disjoint, previously unused
unseen systems, three fresh seeds, clean and noisy conditions, a strong
baseline, and three ablations. Both candidates clear the development screen
and both fail the preregistered unseen criterion. Those negative results are
the observable trigger for a route change, not an inconvenience removed from
the paper.

Second, we evaluate the research-control mechanism itself on a frozen matrix
of \UciTaskCount{} freshly executed UCI tasks and \MDBenchTaskCount{} revealed
MDBench trace-replay tasks. The matrix compares \oneshot{}, \executeonce{},
and the \systemloop{} across \SeedCount{} seeds. Four ablations remove
\vault{} memory, failure feedback, preregistration, or the claim-evidence
gate. The principal endpoint is paired task-success gain over \executeonce{}.
Additional endpoints measure recovery from an initially negative state, exact
reproduction, unsupported scientific claims, research-decision intervention,
wall time, and external cost. Revealed benchmark traces are explicitly
restricted to workflow-behavior evaluation and are not presented as new
holdout evidence for an equation-discovery method.

The study makes three bounded contributions:

\begin{itemize}
  \item It specifies a recursive research state machine in which a scientific
  round is counted only when a parent-linked mechanism change reaches a new,
  precommitted experiment and a visible decision.
  \item It provides an end-to-end audit trail showing two real negative
  method rounds followed by a preregistered systems pivot, with no
  research-decision human intervention after campaign start.
  \item It reports a controlled comparison and component ablations showing
  that the full evidence-bound loop outperforms execute-once on this task
  matrix while preserving exact reruns and zero unsupported claims.
\end{itemize}

The conclusion is deliberately narrow. We do not infer general scientific
autonomy, venue acceptance, or a successful new equation-discovery method.
The experiment tests whether a particular control architecture recovers from
specified workflow failures under a frozen evaluator. Human judgment remains
necessary for novelty, scientific importance, authorship, licensing, venue
fit, and every external submission decision.
